% ============================================================
%  SAS – Příprava k ústní zkoušce B2B37SAS
%  Kompletní vypracování otázek s důrazem na * témata
% ============================================================
\documentclass[12pt,a4paper]{article}

\usepackage[utf8]{inputenc}
\usepackage[T1]{fontenc}
\usepackage[czech]{babel}
\usepackage{lmodern}
\usepackage{amsmath,amssymb,amsthm}
\usepackage{mathtools}
\usepackage{geometry}
\usepackage{graphicx}
\usepackage{xcolor}
\usepackage{tcolorbox}
\usepackage{enumitem}
\usepackage{booktabs}
\usepackage{array}
\usepackage{hyperref}
\usepackage{fancyhdr}
\usepackage{titlesec}
\usepackage{mdframed}
\usepackage{tikz}
\usepackage{pgfplots}
\pgfplotsset{compat=1.18}
\usetikzlibrary{arrows.meta,positioning,shapes}

\geometry{
  top=2.5cm, bottom=2.5cm,
  left=2.5cm, right=2.5cm
}

% ---------- Barvy ----------
\definecolor{hvezdicka}{RGB}{180,0,0}
\definecolor{pozadihv}{RGB}{255,240,240}
\definecolor{pozadiinfo}{RGB}{235,245,255}
\definecolor{pozadimat}{RGB}{240,255,240}
\definecolor{pozadinote}{RGB}{255,253,220}
\definecolor{hlavicka}{RGB}{30,60,120}

% ---------- tcolorbox styly ----------
\tcbuselibrary{skins,breakable}

\newtcolorbox{hvezdBox}[1][]{
  breakable,
  colback=pozadihv,
  colframe=hvezdicka,
  fonttitle=\bfseries\color{white},
  coltitle=white,
  attach boxed title to top left={yshift=-2mm,xshift=4mm},
  boxed title style={colback=hvezdicka,rounded corners},
  title={$\bigstar$ KLÍČOVÁ OTÁZKA},
  #1
}

\newtcolorbox{infoBox}[1][]{
  breakable,
  colback=pozadiinfo,
  colframe=hlavicka!70,
  fonttitle=\bfseries,
  title={#1}
}

\newtcolorbox{matBox}[1][]{
  breakable,
  colback=pozadimat,
  colframe=green!60!black,
  fonttitle=\bfseries,
  title={#1}
}

\newtcolorbox{noteBox}[1][]{
  breakable,
  colback=pozadinote,
  colframe=orange!70!black,
  fonttitle=\bfseries,
  title={#1}
}

% ---------- Záhlaví / zápatí ----------
\pagestyle{fancy}
\fancyhf{}
\fancyhead[L]{\small\textcolor{hlavicka}{\textbf{B2B37SAS} – Příprava k ústní zkoušce}}
\fancyhead[R]{\small\textcolor{hlavicka}{\leftmark}}
\fancyfoot[C]{\thepage}
\renewcommand{\headrulewidth}{0.4pt}

% ---------- Titulky sekcí ----------
\titleformat{\section}[block]
  {\Large\bfseries\color{hlavicka}}
  {\thesection.}{0.6em}{}[\titlerule]

\titleformat{\subsection}[runin]
  {\normalsize\bfseries\color{hvezdicka!80!black}}
  {}{0em}{}[.]

% ---------- Zkratky ----------
\newcommand{\FT}{\mathcal{F}}
\newcommand{\LT}{\mathcal{L}}
\newcommand{\ZT}{\mathcal{Z}}
\newcommand{\Av}{\operatorname{Av}}
\newcommand{\sinc}{\operatorname{sinc}}
\newcommand{\rect}{\operatorname{rect}}
\newcommand{\dirac}{\delta}
\newcommand{\jj}{j}   % imaginární jednotka
\newcommand{\dd}{\mathrm{d}}

\hypersetup{
  colorlinks=true,
  linkcolor=hlavicka,
  urlcolor=blue,
  pdftitle={SAS -- Příprava k ústní zkoušce},
  pdfauthor={B2B37SAS}
}

% ============================================================
\begin{document}

% ---- Titulní strana ----
\begin{titlepage}
  \centering
  \vspace*{2cm}
  {\Huge\bfseries\color{hlavicka} Příprava k ústní zkoušce\\[0.4em]
   \LARGE B2B37SAS – Signály a soustavy}\\[1.5cm]
  {\large Kompletní vypracování okruhů dle BSZZ\\
   se zvláštním důrazem na $\bigstar$ klíčová témata}\\[1cm]
  \begin{tcolorbox}[colback=pozadihv,colframe=hvezdicka,width=0.85\linewidth]
    \centering
    \textbf{Upozornění:} Hvězdičkou ($\bigstar$) jsou označena naprosto základní nosná témata.
    Jejich zásadní neznalost může vést k udělení F \emph{bez ohledu} na předchozí bodové hodnocení.
  \end{tcolorbox}
  \vfill
  {\large\today}
\end{titlepage}

% ---- Obsah ----
\tableofcontents
\newpage

% ============================================================
\section{Klasifikace signálů ve spojitém a diskrétním čase, speciální signály}
% ============================================================

% -----------------------------------------------------------
\begin{hvezdBox}
\subsection*{Co je to signál?}
\end{hvezdBox}

Signál je \textbf{reálná nebo komplexní funkce nezávislé proměnné} (nejčastěji čas $t$ nebo index vzorku $k$). Formálně:

\begin{itemize}[leftmargin=1.5em]
  \item Jakákoliv závislost fyzikální veličiny na čase: $s = s(t)$ nebo $s = s[k]$.
  \item U obrazů nezávislou proměnnou jsou souřadnice $x, y$.
  \item Signál může být popsán matematickou formulací, grafem nebo tabulkou.
\end{itemize}

\begin{noteBox}[Intuice]
Signál = nosič informace. Každý měřitelný fyzikální děj (zvuk, napětí, teplota, elektromagnetická vlna) může být modelován jako signál.
\end{noteBox}

% -----------------------------------------------------------
\begin{hvezdBox}
\subsection*{Jak vypadá signál ve spojitém a diskrétním čase?}
\end{hvezdBox}

\begin{infoBox}[Spojitý čas]
Nezávislá proměnná $t \in \mathbb{R}$ nabývá libovolných reálných hodnot. Signál $s(t)$ je definován pro každý okamžik $t$.
\[
  s: \mathbb{R} \to \mathbb{R} \quad \text{(nebo } \mathbb{C}\text{)}
\]
Příklady: analogový zvukový signál, napětí v elektrickém obvodu.
\end{infoBox}

\begin{infoBox}[Diskrétní čas]
Nezávislá proměnná je celé číslo $k \in \mathbb{Z}$. Signál $s[k]$ je definován pouze pro celočíselné hodnoty $k$ – tzv.\ \textbf{vzorky}.
\[
  s: \mathbb{Z} \to \mathbb{R} \quad \text{(nebo } \mathbb{C}\text{)}
\]
Vzniká vzorkováním spojitého signálu: $s[k] = s(kT_{sa})$, kde $T_{sa}$ je vzorkovací perioda.
\end{infoBox}

% -----------------------------------------------------------
\subsection*{Jak se liší deterministické a náhodné (stochastické) signály?}

\begin{center}
\begin{tabular}{lll}
\toprule
 & \textbf{Deterministický} & \textbf{Náhodný (stochastický)} \\
\midrule
Popis & explicitní matematická funkce & pravděpodobnostní charakteristiky \\
Hodnota v čase $t$ & přesně dána & není předem určena \\
Příklad & $\cos(\omega_0 t)$ & tepelný šum, řeč \\
\bottomrule
\end{tabular}
\end{center}

% -----------------------------------------------------------
\subsection*{Co to jsou kauzální a finitní signály?}

\begin{itemize}
  \item \textbf{Kauzální signál}: $s(t) = 0$ pro $t < 0$. Má definovaný začátek – „příčinnost" (kauzalita).
  \item \textbf{Finitní signál}: má definovaný jak začátek, tak konec – $s(t) = 0$ mimo interval $[t_1, t_2]$. Konečná délka trvání. Každý finitní signál je i kauzální (po posunu).
\end{itemize}

% -----------------------------------------------------------
\begin{hvezdBox}
\subsection*{Speciální signály – vlastnosti}
Jednotkový skok, obdélníkový impuls, jednotkový impuls, Diracův impuls, vzorkovací funkce
\end{hvezdBox}

\subsubsection*{Jednotkový skok $1(t)$, resp.\ $u[k]$}

\begin{matBox}[Definice]
\[
  1(t) = \begin{cases} 0 & t < 0 \\ \tfrac{1}{2} & t = 0 \\ 1 & t > 0 \end{cases}
  \qquad
  u[k] = \begin{cases} 0 & k < 0 \\ 1 & k \geq 0 \end{cases}
\]
\end{matBox}

Vlastnosti:
\begin{itemize}
  \item \textbf{Stejnosměrná složka} (průměrná hodnota): $\Av[1(t)] = \tfrac{1}{2}$.
  \item \textbf{Výkon}: $P = \tfrac{1}{2}$.
  \item \textbf{Autokorelační funkce}: $R(\tau) = \tfrac{1}{2}$ (konstantní).
  \item Výkonový signál (nekonečná energie, konečný výkon).
\end{itemize}

\subsubsection*{Obdélníkový impuls $\rect(t)$}

\begin{matBox}[Definice]
\[
  \rect\!\left(\frac{t}{\tau}\right) = \begin{cases} 1 & |t| < \tau/2 \\ 0 & |t| > \tau/2 \end{cases}
\]
\end{matBox}

Složen ze dvou jednotkových skoků: $\rect\!\left(\tfrac{t}{\tau}\right) = 1\!\left(t + \tfrac{\tau}{2}\right) - 1\!\left(t - \tfrac{\tau}{2}\right)$.\\
Energetický signál. Jeho Fourierova transformace je vzorkovací funkce (sinc):
\[
  \FT\{\rect(t/\tau)\} = \tau \cdot \sinc\!\left(\frac{\omega\tau}{2\pi}\right).
\]

\subsubsection*{Jednotkový impuls $\delta[k]$ (Kroneckerovo delta)}

\begin{matBox}[Definice]
\[
  \delta[k] = \begin{cases} 1 & k = 0 \\ 0 & k \neq 0 \end{cases}
\]
\end{matBox}

Jde o Diracův impuls v \emph{diskrétním} čase. Má vzorkovací vlastnost: $s[k] * \delta[k] = s[k]$.

\subsubsection*{Diracův impuls $\delta(t)$}

\begin{matBox}[Definice (distribuce)]
\[
  \delta(t) = 0 \;\text{ pro }\; t \neq 0, \qquad \int_{-\infty}^{+\infty} \delta(t)\,\dd t = 1.
\]
\end{matBox}

Klíčové vlastnosti:
\begin{itemize}
  \item \textbf{Derivace} jednotkového skoku: $\delta(t) = \dfrac{\dd}{\dd t}1(t)$.
  \item \textbf{Jednotková plocha} pod impulsem.
  \item \textbf{Vzorkovací (sifting) vlastnost}:
    \[
      \int_{-\infty}^{+\infty} s(t)\,\delta(t - t_0)\,\dd t = s(t_0).
    \]
  \item Fourierova transformace: $\FT\{\delta(t)\} = 1$ (konstantní spektrum).
\end{itemize}

\subsubsection*{Vzorkovací funkce (sinc)}

\begin{matBox}[Definice]
\[
  \operatorname{Sa}(t) = \frac{\sin t}{t}, \quad \operatorname{Sa}(0) = 1
  \qquad
  \sinc(t) = \frac{\sin(\pi t)}{\pi t} \quad \text{(normalizovaná)}
\]
\end{matBox}

Vlastnosti:
\begin{itemize}
  \item Energetický signál: $E = \pi$ (pro $\operatorname{Sa}$) nebo $E=1$ (pro $\sinc$).
  \item Nuly v $t = n\pi$, $n \in \mathbb{Z} \setminus \{0\}$.
  \item Fourierova transformace: $\FT\{\sinc(t/\tau)\} = \tau\cdot\rect(\omega\tau/2\pi)$ – \emph{obdélník ve spektru}.
\end{itemize}

% -----------------------------------------------------------
\subsection*{Vzorkovací vlastnost Diracova a jednotkového impulsu}

Pronásobení libovolné funkce $s(t)$ Diracem $\delta(t-t_0)$ vyselektuje hodnotu funkce v bodě $t_0$:
\[
  s(t)\cdot\delta(t-t_0) = s(t_0)\cdot\delta(t-t_0).
\]
Po integraci pak:
\[
  \int_{-\infty}^{+\infty} s(t)\,\delta(t-t_0)\,\dd t = s(t_0).
\]
Analogicky v diskrétním čase: $\sum_{k=-\infty}^{\infty} s[k]\,\delta[k-k_0] = s[k_0]$.

\newpage
% ============================================================
\section{Časová a spektrální reprezentace signálů, korelace, základní teorémy}
% ============================================================

\subsection*{Energie a výkon signálu}

\begin{matBox}[Energie a výkon – spojitý čas]
\[
  E = \int_{-\infty}^{+\infty} |s(t)|^2\,\dd t
  \qquad
  P = \Av\!\left[|s(t)|^2\right] = \lim_{T\to\infty}\frac{1}{2T}\int_{-T}^{T}|s(t)|^2\,\dd t
\]
\end{matBox}

\begin{matBox}[Energie a výkon – diskrétní čas]
\[
  E = \sum_{k=-\infty}^{\infty} |s[k]|^2
  \qquad
  P = \lim_{N\to\infty}\frac{1}{2N+1}\sum_{k=-N}^{N}|s[k]|^2
\]
\end{matBox}

\begin{center}
\begin{tabular}{lll}
\toprule
 & \textbf{Energetický signál} & \textbf{Výkonový signál} \\
\midrule
Energie $E$ & konečná ($E < \infty$) & nekonečná \\
Výkon $P$ & nulový & konečný a nenulový \\
Příklady & finitní signály, $\sinc$ & periodické signály, $1(t)$ \\
\bottomrule
\end{tabular}
\end{center}

\begin{noteBox}[Důležité]
Pokud je signál finitní $\Rightarrow$ je energetický (ale ne naopak). Periodický signál $\Rightarrow$ výkonový.
\end{noteBox}

% -----------------------------------------------------------
\begin{hvezdBox}
\subsection*{Vzájemná korelační funkce a autokorelační funkce}
\end{hvezdBox}

\subsubsection*{Vzájemná korelační funkce $R_{12}(\tau)$}

Míra \textbf{podobnosti dvou různých signálů} $s_1, s_2$ s uvažováním vzájemného posunu $\tau$:

\begin{matBox}[Vzorec]
\[
  R_{12}(\tau) = \int_{-\infty}^{+\infty} s_1(t)\,s_2^*(t-\tau)\,\dd t
  \qquad
  R_{12}[m] = \sum_{k=-\infty}^{\infty} s_1[k]\,s_2^*[k-m]
\]
\end{matBox}

Mechanismus: posouváme $s_2$ o $\tau$, pronásobíme s $s_1$ a integrujeme (sečteme) – hledáme posun $\tau$, při němž jsou si signály nejpodobnější.

Pro reálné signály: $R_{12}(\tau) = R_{21}(-\tau)$.

\subsubsection*{Autokorelační funkce $R(\tau)$}

Míra \textbf{závislosti signálu na sobě samém} s posunem $\tau$:

\begin{matBox}[Vzorec – energetický signál]
\[
  R(\tau) = \int_{-\infty}^{+\infty} s(t)\,s^*(t-\tau)\,\dd t
\]
\end{matBox}

\begin{matBox}[Vzorec – výkonový signál]
\[
  R(\tau) = \Av\!\left[s(t+\tau)\,s^*(t)\right]
\]
\end{matBox}

Klíčové vlastnosti ACF:
\begin{itemize}
  \item \textbf{Sudá funkce}: $R(\tau) = R(-\tau)$ (pro reálné signály).
  \item \textbf{Globální maximum v $\tau = 0$}: $|R(\tau)| \leq R(0)$.
  \item $R(0) = E$ (energetické signály) nebo $R(0) = P$ (výkonové signály).
\end{itemize}

Využití: GPS (korelace přijatého signálu s referenčními kódy), detekce opakujících se vzorů.

% -----------------------------------------------------------
\subsection*{Vztah ACF a energie/výkonu}

\begin{matBox}[]
\[
  R(0) = E \quad \text{(energetický signál)}
  \qquad
  R(0) = P \quad \text{(výkonový signál)}
\]
\end{matBox}

% -----------------------------------------------------------
\subsection*{Operátor $\Av$ a stejnosměrná složka}

\[
  s_{SS} = \Av[s(t)] = \lim_{T\to\infty}\frac{1}{2T}\int_{-T}^{T}s(t)\,\dd t
\]

Stejnosměrná složka odpovídá koeficientu $c_0$ Fourierovy řady (nultá harmonická).

% -----------------------------------------------------------
\begin{hvezdBox}
\subsection*{Ortogonalita dvou signálů}
\end{hvezdBox}

Dva signály $s_1(t)$ a $s_2(t)$ jsou \textbf{vzájemně ortogonální}, pokud jejich vzájemná energie (resp.\ výkon) je nulová:

\begin{matBox}[Podmínka ortogonality]
\[
  E_{12} = \int_{-\infty}^{+\infty} s_1(t)\,s_2^*(t)\,\dd t = 0
  \quad \Leftrightarrow \quad
  s_1 \perp s_2
\]
\end{matBox}

\begin{itemize}
  \item Ortogonální signály se \textbf{vzájemně neovlivňují} – jejich energetický „průnik" je nulový.
  \item Příklady: $\sin(n\omega_0 t)$ a $\cos(m\omega_0 t)$ pro $n \neq m$, nebo $\sin(n\omega_0 t)$ a $\sin(m\omega_0 t)$ pro $n \neq m$ (základ Fourierovy řady).
  \item Ortogonalita je nutnou podmínkou pro jednoznačný rozklad signálu do bázových funkcí.
\end{itemize}

% -----------------------------------------------------------
\begin{hvezdBox}
\subsection*{Rozklad periodického signálu do Fourierovy řady (FS)}
\end{hvezdBox}

Každý periodický signál $s(t)$ s periodou $T_0$ lze rozložit do \textbf{nekonečné lineární kombinace komplexních harmonických funkcí}:

\begin{matBox}[Fourierova řada – syntéza (zpětná transformace)]
\[
  s(t) = \sum_{n=-\infty}^{+\infty} c_n\, e^{j n \omega_0 t}, \qquad \omega_0 = \frac{2\pi}{T_0}
\]
\end{matBox}

\begin{matBox}[Výpočet koeficientů – analýza]
\[
  c_n = \frac{1}{T_0}\int_{T_0} s(t)\, e^{-j n \omega_0 t}\,\dd t
\]
\end{matBox}

Princip:
\begin{enumerate}
  \item Bázové funkce $\{e^{jn\omega_0 t}\}$ tvoří \textbf{úplný ortogonální systém} – vzájemně se neovlivňují.
  \item Koeficient $c_n$ vyjadřuje, „jakou měrou se $n$-tá harmonická vyskytuje v signálu $s(t)$".
  \item Dělení $1/T_0$ zajišťuje ortonormalitu systému.
  \item Výsledné spektrum je \textbf{diskrétní, neperiodické} (čárové spektrum).
\end{enumerate}

\begin{noteBox}[Výkonové spektrum]
Výkon $n$-té harmonické odpovídá $|c_n|^2$. Celkový výkon signálu:
\[
  P = \sum_{n=-\infty}^{+\infty} |c_n|^2 = R(0)
\]
(Parsevalova věta pro FS)
\end{noteBox}

% -----------------------------------------------------------
\begin{hvezdBox}
\subsection*{Fourierova transformace (FT) – spektrum obecných signálů}
\end{hvezdBox}

FT zobecňuje principy FS na \textbf{neperiodické spojité signály} (limitní přechod $T_0 \to \infty$, $\omega_0 \to 0$):

\begin{matBox}[Fourierova transformace – přímá]
\[
  S(\omega) = \FT\{s(t)\} = \int_{-\infty}^{+\infty} s(t)\, e^{-j\omega t}\,\dd t
\]
\end{matBox}

\begin{matBox}[Zpětná Fourierova transformace]
\[
  s(t) = \FT^{-1}\{S(\omega)\} = \frac{1}{2\pi}\int_{-\infty}^{+\infty} S(\omega)\, e^{j\omega t}\,\dd\omega
\]
\end{matBox}

Výsledné spektrum je \textbf{spojité, neperiodické}.

Vlastnosti FT reálného signálu ($s(t) \in \mathbb{R}$):
\[
  |S(\omega)| = |S(-\omega)| \quad \text{(sudá amplitudová charakteristika)}
  \qquad
  \arg S(\omega) = -\arg S(-\omega) \quad \text{(lichá fázová charakteristika)}
\]

% -----------------------------------------------------------
\begin{hvezdBox}
\subsection*{Přehled spekter základních signálů}
\end{hvezdBox}

\begin{center}
\begin{tabular}{lll}
\toprule
\textbf{Signál v čase} & \textbf{Spektrum} $S(\omega)$ & \textbf{Poznámka} \\
\midrule
$\delta(t)$ & $1$ & konstanta \\
$1$ (konstantní) & $2\pi\,\delta(\omega)$ & Dirac na $\omega=0$ \\
$\operatorname{Sa}(\omega_0 t)$ & $\dfrac{\pi}{\omega_0}\,\rect\!\left(\dfrac{\omega}{2\omega_0}\right)$ & obdélník \\
$\rect(t/\tau)$ & $\tau\,\operatorname{Sa}(\omega\tau/2)$ & Sa funkce \\
$e^{j\omega_0 t}$ & $2\pi\,\delta(\omega-\omega_0)$ & Dirac v $\omega_0$ \\
$\cos(\omega_0 t)$ & $\pi[\delta(\omega-\omega_0)+\delta(\omega+\omega_0)]$ & dva Diraci \\
$\sin(\omega_0 t)$ & $\dfrac{\pi}{j}[\delta(\omega-\omega_0)-\delta(\omega+\omega_0)]$ & antisymetricky \\
\bottomrule
\end{tabular}
\end{center}

\begin{noteBox}[Dualita FT]
FT je \textbf{duální}: obdélník $\leftrightarrow$ sinc, Dirac $\leftrightarrow$ konstanta. Spektrum konstantního signálu je koncentrováno do jediného bodu (Dirac), a naopak.
\end{noteBox}

% -----------------------------------------------------------
\begin{hvezdBox}
\subsection*{Věta o FT konvoluce}
\end{hvezdBox}

\begin{matBox}[Konvoluční věta]
\[
  \FT\{s_1(t) * s_2(t)\} = S_1(\omega)\cdot S_2(\omega)
\]
\[
  \FT\{s_1(t)\cdot s_2(t)\} = \frac{1}{2\pi}\,S_1(\omega) * S_2(\omega)
\]
\end{matBox}

\textbf{Fyzikální význam}: konvoluce v časové oblasti odpovídá \emph{násobení} ve spektrální oblasti a naopak. To je základ pro:
\begin{itemize}
  \item Popis LTI soustav: $Y(\omega) = H(\omega)\cdot X(\omega)$.
  \item Filtraci: ve spektrální oblasti stačí pronásobit; v časové oblasti bychom museli konvolvovat.
  \item Vzorkování: vzorkování = násobení Diracovým hřebenem v čase = konvoluce (periodizace) ve spektru.
\end{itemize}

% -----------------------------------------------------------
\subsection*{Přehled Fourierovských transformací}

\begin{center}
\begin{tabular}{lllll}
\toprule
\textbf{Transformace} & \textbf{Signál} & \textbf{Spektrum} & & \\
 & (čas) & (frekvence) & & \\
\midrule
FS  & spojitý, \emph{periodický} & diskrétní, neperiodické & & \\
DFS & diskrétní, \emph{periodický} & diskrétní, periodické & & \\
FT  & spojitý, \emph{neperiodický} & spojité, neperiodické & & \\
DtFT& diskrétní, \emph{neperiodický} & spojité, \emph{periodické} & & \\
DFT & diskrétní, \emph{finitní} ($N$ vzorků) & diskrétní, periodické ($N$ hodnot) & & \\
\bottomrule
\end{tabular}
\end{center}

% -----------------------------------------------------------
\begin{hvezdBox}
\subsection*{DtFT – základní vlastnost spektra}
\end{hvezdBox}

Diskrétní Fourierova transformace v čase (DtFT) zobrazuje diskrétní neperiodický signál na \textbf{spojité, periodické spektrum}:

\begin{matBox}[DtFT]
\[
  S(e^{j\Omega}) = \sum_{k=-\infty}^{\infty} s[k]\, e^{-j\Omega k}
\]
\end{matBox}

Spektrum je \textbf{periodické s periodou $2\pi$} (nebo $\omega_{sa}$ v závislosti na normování).

\begin{noteBox}[Proč periodické?]
Diskrétní signál vznikl vzorkováním – ve spektru se proto objevují kopie spektra s periodou $\omega_{sa}$. DtFT tuto periodicitu explicitně zahrnuje.
\end{noteBox}

% -----------------------------------------------------------
\subsection*{DFT a FFT}

\begin{matBox}[DFT – vztahy]
\[
  D[n] = \sum_{k=0}^{N-1} d[k]\, e^{-jnk\frac{2\pi}{N}}
  \qquad
  d[k] = \frac{1}{N}\sum_{n=0}^{N-1} D[n]\, e^{jnk\frac{2\pi}{N}}
\]
\end{matBox}

\begin{itemize}
  \item \textbf{DFT}: umožňuje numerický výpočet spektra z konečného počtu vzorků. Pracuje jen s $N$ vzorky – nepotřebujeme znát signál na celém definičním oboru.
  \item \textbf{FFT}: efektivní algoritmus výpočtu DFT, snižuje počet operací z $O(N^2)$ na $O(N\log_2 N)$.
\end{itemize}

\newpage
% ============================================================
\section{Vzorkování a interpolace signálu}
% ============================================================

\begin{hvezdBox}
\subsection*{Shannonův vzorkovací teorém}
\end{hvezdBox}

\textbf{Formulace}: Spektrálně omezený signál $s(t)$ s maximálním kmitočtem $f_m$ lze \emph{jednoznačně rekonstruovat} ze svých vzorků, pokud je vzorkovací kmitočet $f_{sa}$ dostatečně vysoký:

\begin{matBox}[Shannonova (Nyquistova) podmínka]
\[
  f_{sa} > 2 f_m
  \qquad \Leftrightarrow \qquad
  \omega_{sa} > 2\omega_m
\]
\end{matBox}

\subsubsection*{Odvození ve spektrální oblasti}

\begin{enumerate}
  \item \textbf{Ideální vzorkování} = násobení signálu Diracovým hřebenem $p(t) = \sum_k \delta(t - kT_{sa})$:
    \[
      s_a(t) = s(t)\cdot p(t) = \sum_{k=-\infty}^{\infty} s(kT_{sa})\,\delta(t-kT_{sa})
    \]
  \item Ve \textbf{spektrální oblasti}: FT Diracového hřebene je opět hřeben s periodou $\omega_{sa}$, takže:
    \[
      S_a(\omega) = \frac{1}{T_{sa}}\sum_{n=-\infty}^{\infty} S(\omega - n\omega_{sa})
    \]
    Spektrum vzorkovaného signálu je \textbf{periodická kopie} originálu s periodou $\omega_{sa}$.
  \item \textbf{Rekonstrukce}: pokud jsou kopie spektra \emph{odděleny} (nepřekrývají se), lze je ideální dolní propustí oddělit a dostat zpět $S(\omega)$.
  \item Kopie se nepřekrývají $\Leftrightarrow$ $\omega_{sa} > 2\omega_m$.
\end{enumerate}

\begin{noteBox}[Aliasing]
Pokud $f_{sa} \leq 2f_m$, dochází k \textbf{aliasingu} – překryvu spektrálních kopií. Rekonstrukce je zkreslená a jednoznačná obnova signálu není možná.
\end{noteBox}

% -----------------------------------------------------------
\begin{hvezdBox}
\subsection*{Ideální interpolace vzorkovaného signálu}
\end{hvezdBox}

Rekonstrukce (interpolace) spojitého signálu ze vzorků $s[k] = s(kT_{sa})$:

\begin{matBox}[Ideální interpolace – Whittaker–Shannon]
\[
  s(t) = \sum_{k=-\infty}^{\infty} s[k]\cdot\sinc\!\left(\frac{\omega_{sa}}{2}(t - kT_{sa})\right)
       = \sum_{k=-\infty}^{\infty} s(kT_{sa})\cdot\operatorname{Sa}\!\left(\frac{\pi(t-kT_{sa})}{T_{sa}}\right)
\]
\end{matBox}

Postup:
\begin{enumerate}
  \item Vzorky $s[k]$ nahradíme váhovanými Dirakovými impulsy (ideálně vzorkovaný signál).
  \item Výsledek váhujeme \textbf{ideální dolní propustí} s mezním kmitočtem $\omega_{sa}/2$.
  \item DP má v časové oblasti impulsovou odezvu = $\sinc$ funkce.
\end{enumerate}

\subsection*{Problémy realizace ideálního interpolátoru}

\begin{itemize}
  \item Ideální dolní propust je \textbf{nekauzální} (impulsová odezva sahá do záporných časů) – fyzikálně nerealizovatelná.
  \item Diracovy impulsy nelze v praxi generovat.
\end{itemize}

\textbf{Řešení v reálných aplikacích}:
\begin{itemize}
  \item Vzorky nahradíme obdélníkovými pulzy (zero-order hold nebo linear interpolation).
  \item Dolní propust realizujeme kauzálním FIR/IIR filtrem (přijatelné zkreslení).
\end{itemize}

\newpage
% ============================================================
\section{Klasifikace soustav, LTI soustavy, konvoluce, stabilita}
% ============================================================

\begin{hvezdBox}
\subsection*{Co je to soustava a operátor soustavy?}
\end{hvezdBox}

\textbf{Soustava} (systém) = objekt, který pod vlivem \emph{buzení} (vstupního signálu $x$) produkuje \emph{odezvu} (výstupní signál $y$):
\[
  y = \mathcal{A}\{x\}
\]

$\mathcal{A}$ je \textbf{operátor soustavy}. Soustava dává do vztahu skupinu signálů (fyzikálních nebo matematických veličin).

% -----------------------------------------------------------
\subsection*{Klasifikace soustav}

\begin{center}
\begin{tabular}{ll}
\toprule
\textbf{Vlastnost} & \textbf{Popis} \\
\midrule
Kauzální & $y(t)$ závisí jen na $x(\tau)$, $\tau \leq t$ (ne na budoucnosti) \\
Nekauzální & $y(t)$ závisí i na budoucích hodnotách vstupu \\
Bez paměti & $y(t)$ závisí jen na $x(t)$ (aktuální vstup) \\
S pamětí & $y(t)$ závisí na minulých nebo budoucích hodnotách \\
Stabilní (BIBO) & omezený vstup $\Rightarrow$ omezený výstup \\
\bottomrule
\end{tabular}
\end{center}

\subsection*{BIBO stabilita}

\textbf{Bounded-Input Bounded-Output}: Soustava je stabilní, pokud pro každý omezený vstup ($|x(t)| \leq M_x < \infty$) je výstup také omezený ($|y(t)| \leq M_y < \infty$).

% -----------------------------------------------------------
\begin{hvezdBox}
\subsection*{LTI soustavy (Linear Time-Invariant)}
\end{hvezdBox}

Soustava $\mathcal{A}$ je LTI, pokud splňuje:

\begin{enumerate}
  \item \textbf{Linearita} (princip superpozice):
    \[
      \mathcal{A}\{\alpha x_1(t) + \beta x_2(t)\} = \alpha\,\mathcal{A}\{x_1(t)\} + \beta\,\mathcal{A}\{x_2(t)\}
    \]
  \item \textbf{Časová invariance}: posun vstupu o $t_0$ způsobí stejný posun výstupu:
    \[
      \mathcal{A}\{x(t-t_0)\} = y(t-t_0)
    \]
\end{enumerate}

LTI soustavy jsou \textbf{klíčové}, protože jsou plně popsány jedinou funkcí – impulsovou odezvou $h(t)$.

% -----------------------------------------------------------
\begin{hvezdBox}
\subsection*{Impulsová odezva LTI soustavy}
\end{hvezdBox}

\begin{matBox}[Definice]
\[
  h(t) = \mathcal{A}\{\delta(t)\} \quad \text{(spojitý čas)}
  \qquad
  h[k] = \mathcal{A}\{\delta[k]\} \quad \text{(diskrétní čas)}
\]
\end{matBox}

Impulsová odezva $h(t)$ \textbf{plně charakterizuje} LTI soustavu:
\begin{itemize}
  \item Kauzalita: $h(t) = 0$ pro $t < 0$.
  \item Stabilita: $\displaystyle\int_{-\infty}^{+\infty} |h(t)|\,\dd t < \infty$ (absolutní integrabilita).
\end{itemize}

% -----------------------------------------------------------
\begin{hvezdBox}
\subsection*{Konvoluce – princip a výpočet}
\end{hvezdBox}

Výstup LTI soustavy je \textbf{konvoluce} vstupu $x$ s impulsovou odezvou $h$:

\begin{matBox}[Konvoluce – spojitý čas]
\[
  y(t) = x(t) * h(t) = \int_{-\infty}^{+\infty} x(\tau)\,h(t-\tau)\,\dd\tau
\]
\end{matBox}

\begin{matBox}[Konvoluce – diskrétní čas]
\[
  y[k] = x[k] * h[k] = \sum_{i=-\infty}^{\infty} x[i]\,h[k-i]
\]
\end{matBox}

\textbf{Grafický postup výpočtu konvoluce}:
\begin{enumerate}
  \item $s_1(\tau)$ – první signál jako funkce proměnné $\tau$.
  \item $s_2(-\tau)$ – druhý signál \emph{převrácený} (zrcadlení podle osy).
  \item $s_2(t-\tau)$ – převrácený signál \emph{posunutý} o $t$ doprava.
  \item Integrál (plocha pod grafem) součinu $s_1(\tau)\cdot s_2(t-\tau)$ = hodnota konvoluce pro daný $t$.
  \item Opakujeme pro všechna $t$.
\end{enumerate}

Vlastnosti konvoluce: komutativní ($a*b = b*a$), asociativní, distributivní.

% -----------------------------------------------------------
\begin{hvezdBox}
\subsection*{Kauzalita a stabilita z impulsové odezvy}
\end{hvezdBox}

\begin{infoBox}[Kauzalita]
Soustava je kauzální $\Leftrightarrow$ $h(t) = 0$ pro $t < 0$ (resp.\ $h[k] = 0$ pro $k < 0$).
\end{infoBox}

\begin{infoBox}[BIBO stabilita]
Soustava je BIBO stabilní $\Leftrightarrow$ impulsová odezva je absolutně integrabilní (sumovatelná):
\[
  \int_{-\infty}^{+\infty} |h(t)|\,\dd t < \infty
  \qquad
  \sum_{k=-\infty}^{\infty} |h[k]| < \infty
\]
\end{infoBox}

% -----------------------------------------------------------
\begin{hvezdBox}
\subsection*{Systémová funkce LTI soustavy – nuly a póly}
\end{hvezdBox}

\textbf{Systémová funkce} $H(s)$ (resp.\ $H(z)$) je Laplaceova (resp.\ Z-) transformace impulsové odezvy:

\begin{matBox}[Systémová funkce]
\[
  H(s) = \LT\{h(t)\} \quad \text{(spojitý čas)}
  \qquad
  H(z) = \ZT\{h[k]\} \quad \text{(diskrétní čas)}
\]
\end{matBox}

Pro LTI soustavu popsanou diferenciální rovnicí s konstantními koeficienty je $H(s)$ \textbf{racionální lomená funkce}:
\[
  H(s) = \frac{B(s)}{A(s)} = \frac{b_M s^M + \cdots + b_0}{a_N s^N + \cdots + a_0}
\]

\begin{itemize}
  \item \textbf{Nuly}: kořeny čitatele $B(s) = 0$ – kmitočty, kde přenos je nulový.
  \item \textbf{Póly}: kořeny jmenovatele $A(s) = 0$ – charakteristické frekvence soustavy.
\end{itemize}

\begin{infoBox}[Stabilita z polohy pólů]
Spojitá LTI soustava je BIBO stabilní $\Leftrightarrow$ všechny póly leží v \textbf{levé polorovině} komplexní roviny ($\operatorname{Re}(p_i) < 0$).\\[4pt]
Diskrétní LTI soustava je BIBO stabilní $\Leftrightarrow$ všechny póly leží \textbf{uvnitř jednotkové kružnice} ($|p_i| < 1$).
\end{infoBox}

% -----------------------------------------------------------
\begin{hvezdBox}
\subsection*{Frekvenční charakteristika LTI soustavy}
\end{hvezdBox}

\textbf{Přenosová funkce (frekvenční charakteristika)} $H(\omega)$: popisuje, jak soustava zpracovává harmonické signály.

\begin{matBox}[Vztah]
\[
  H(\omega) = H(s)\big|_{s=j\omega} = \FT\{h(t)\}
  \quad \text{(spojitá soustava)}
\]
\[
  H(e^{j\Omega}) = H(z)\big|_{z=e^{j\Omega}} = \text{DtFT}\{h[k]\}
  \quad \text{(diskrétní soustava)}
\]
\end{matBox}

\begin{itemize}
  \item \textbf{Amplitudová charakteristika} $|H(\omega)|$: jak soustava mění amplitudu na kmitočtu $\omega$.
  \item \textbf{Fázová charakteristika} $\arg H(\omega)$: fázový posun.
  \item Pro reálnou impulsovou odezvu: $|H(\omega)|$ je sudá, $\arg H(\omega)$ je lichá.
  \item Frekvenční charakteristika diskrétní soustavy je \textbf{periodická} s periodou $2\pi$ (resp.\ $\omega_{sa}$).
\end{itemize}

Odezva LTI soustavy na harmonické buzení $x(t) = A\cos(\omega_0 t + \varphi)$:
\[
  y(t) = A\cdot|H(\omega_0)|\cdot\cos\!\bigl(\omega_0 t + \varphi + \arg H(\omega_0)\bigr)
\]
LTI soustava mění jen \emph{amplitudu} a \emph{fázi}, kmitočet zůstává stejný.

\subsection*{Membránový model}

Vizualizace systémové funkce $H(s)$ v komplexní rovině: přišpendlení nul (přenos = 0), vytáhnutí pólů (přenos $\to \infty$). Přenosová funkce $|H(j\omega)|$ odpovídá výšce „membrány" podél imaginární osy.

\newpage
% ============================================================
\section{Pásmové signály, komplexní obálka, Hilbertova transformace}
% ============================================================

\begin{hvezdBox}
\subsection*{Reálný pásmový signál – vlastnosti}
\end{hvezdBox}

Pásmový signál (Band-Pass, BP) má spektrum soustředěné \textbf{okolo nosného kmitočtu $\omega_c$}:

\begin{itemize}
  \item Spektrum je nenulové pouze v pásmech $[\omega_d, \omega_h]$ a $[-\omega_h, -\omega_d]$.
  \item \textbf{Symetrické amplitudové spektrum}: $|S(\omega)| = |S(-\omega)|$ (reálný signál).
  \item Nosný kmitočet $\omega_c \gg$ šířka pásma $B$: $\omega_c \gg 2\pi B$.
  \item Lze zapsat jako: $s(t) = V(t)\cos(\omega_c t + \Theta(t))$, kde $V(t)$ je \emph{obálka} a $\Theta(t)$ je \emph{fáze}.
\end{itemize}

% -----------------------------------------------------------
\begin{hvezdBox}
\subsection*{Přechod: pásmový signál → analytický signál → komplexní obálka}
\end{hvezdBox}

\begin{enumerate}
  \item \textbf{Reálný pásmový signál $s(t)$}:\\
    Symetrické spektrum $S(\omega)$, nenulové okolo $\pm\omega_c$.

  \item \textbf{Analytický signál $z(t)$}:\\
    Odstraníme záporné kmitočty a vynásobíme kladnou část spektra dvěma:
    \[
      Z(\omega) = 2\cdot\mathbf{1}(\omega)\cdot S(\omega) = \begin{cases} 2S(\omega) & \omega > 0 \\ 0 & \omega < 0 \end{cases}
    \]
    Analytický signál je \textbf{komplexní}: $z(t) = s(t) + j\hat{s}(t)$, kde $\hat{s}(t)$ je Hilbertova transformace $s(t)$.

  \item \textbf{Komplexní obálka $\tilde{s}(t)$ (CE – Complex Envelope)}:\\
    Posunutí spektra analytického signálu do nulového kmitočtu (dolů o $\omega_c$):
    \[
      \tilde{S}(\omega) = Z(\omega + \omega_c)
      \qquad \Leftrightarrow \qquad
      \tilde{s}(t) = z(t)\cdot e^{-j\omega_c t}
    \]
\end{enumerate}

Komplexní obálka je \textbf{nízkofrekvenční (NF) komplexní reprezentace} vysokofrekvenčního pásmového signálu.

Zpětný přechod:
\[
  s(t) = \operatorname{Re}\!\left\{\tilde{s}(t)\cdot e^{j\omega_c t}\right\}
\]

\subsubsection*{Soufázová a kvadraturní složka}
\[
  \tilde{s}(t) = a(t) + j\,c(t)
\]
\begin{itemize}
  \item $a(t) = \operatorname{Re}\{\tilde{s}(t)\}$ – \textbf{soufázová} (in-phase, I) složka.
  \item $c(t) = \operatorname{Im}\{\tilde{s}(t)\}$ – \textbf{kvadraturní} (quadrature, Q) složka.
\end{itemize}

Výhody komplexní obálky:
\begin{itemize}
  \item Nižší vzorkovací kmitočet (signál je NF).
  \item Jednodušší matematický popis a simulace.
  \item Snadnější realizace algoritmů zpracování.
\end{itemize}

% -----------------------------------------------------------
\begin{hvezdBox}
\subsection*{Hilbertova transformace}
\end{hvezdBox}

Hilbertova transformace (HT) je LTI soustava definovaná přenosovou funkcí:

\begin{matBox}[Hilbertův transformátor]
\[
  H(\omega) = -j\cdot\operatorname{sgn}(\omega) = \begin{cases} -j & \omega > 0 \\ +j & \omega < 0 \end{cases}
\]
Impulsová odezva: $h(t) = \dfrac{1}{\pi t}$
\end{matBox}

\begin{itemize}
  \item \textbf{Lineární} soustava (konvoluce).
  \item Na \emph{kladných} kmitočtech realizuje \textbf{fázové zpoždění o $\pi/2$}.
  \item Na \emph{záporných} kmitočtech realizuje fázový předstih o $\pi/2$.
  \item Amplitudová charakteristika je konstantní: $|H(\omega)| = 1$.
  \item Soustava je \textbf{nekauzální} ($h(t) \neq 0$ pro $t < 0$) a \textbf{nestabilní} (není absolutně integrabilní).
\end{itemize}

Využití: výpočet analytického signálu $z(t) = s(t) + j\hat{s}(t)$.

\subsection*{Metody přechodu BP $\leftrightarrow$ CE}

\textbf{Metoda fázového posunu}:
\begin{enumerate}
  \item Vypočítáme $\hat{s}(t)$ = Hilbertova transformace $s(t)$.
  \item Analytický signál: $z(t) = s(t) + j\hat{s}(t)$.
  \item Komplexní obálka: $\tilde{s}(t) = z(t)\cdot e^{-j\omega_c t}$.
\end{enumerate}

\textbf{Filtrační metoda}:
\begin{enumerate}
  \item Soufázová složka: $a(t) = \operatorname{LP}\{s(t)\cdot 2\cos(\omega_c t)\}$ (násobení cosinem + DP).
  \item Kvadraturní složka: $c(t) = \operatorname{LP}\{-s(t)\cdot 2\sin(\omega_c t)\}$ (násobení sinem + DP).
\end{enumerate}

\subsection*{Vzorkování pásmových signálů}

Vzorkovací podmínka pro BP signál v pásmu $[\omega_d, \omega_h]$, šířka pásma $B = \omega_h - \omega_d$:
\[
  \frac{2\omega_h}{M+1} \leq \omega_{sa} \leq \frac{2\omega_d}{M}
\]
kde $M$ je \emph{číslo pásma}. Pro $M=0$ (základní pásmo) dostaneme Shannonovu podmínku $\omega_{sa} \geq 2\omega_h$.

\newpage
% ============================================================
\section{Typy základních analogových modulací}
% ============================================================

\begin{hvezdBox}
\subsection*{Základní pojmy modulace}
\end{hvezdBox}

\textbf{Modulace} = ovlivnění parametrů \emph{nosné vlny} (carrier) \emph{modulačním signálem}:
\[
  s(t) = V(t)\cdot\cos\!\bigl(\omega_c t + \Theta(t)\bigr)
\]

\begin{itemize}
  \item \textbf{Nosná vlna} $s_c(t) = A_c\cos(\omega_c t)$: vysokofrekvenční harmonický signál vhodný pro přenos médiem (RF vlna, optický paprsek).
  \item \textbf{Modulační signál} $s_M(t)$: nese informaci (zvuk, data, video).
  \item \textbf{Modulovaný signál} $s(t)$: výsledek modulace, přenáší informaci vhodnou formou pro dané médium.
\end{itemize}

Dělení analogových modulací:
\begin{center}
\begin{tabular}{lll}
\toprule
\textbf{Skupina} & \textbf{Co modulační signál ovlivňuje} & \textbf{Typ} \\
\midrule
Amplitudová & amplitudu $V(t)$, fáze $\Theta$ = konst. & AM-DSB, AM-SSB, AM-VSB \\
Úhlová & fázi/kmitočet $\Theta(t)$, obálka $V$ = konst. & PM, FM \\
\bottomrule
\end{tabular}
\end{center}

% -----------------------------------------------------------
\begin{hvezdBox}
\subsection*{Amplitudová modulace AM (Amplitude Modulation)}
\end{hvezdBox}

\begin{matBox}[AM-DSB s nosnou vlnou]
\[
  s_{AM}(t) = A_c\bigl[1 + m\cdot s_M(t)\bigr]\cos(\omega_c t)
\]
kde $m$ je \emph{modulační index} (hloubka modulace).
\end{matBox}

\textbf{Princip v časové oblasti}: obálka nosné vlny kopíruje tvar modulačního signálu $s_M(t)$.

\textbf{Princip ve spektrální oblasti}:
\begin{itemize}
  \item Spektrum modulačního signálu $S_M(\omega)$ se přesune na kmitočty $\pm\omega_c$.
  \item Výsledek: horní postranní pásmo USB ($\omega > \omega_c$) + dolní postranní pásmo LSB ($\omega < \omega_c$) + nosná.
\end{itemize}

\textbf{Nevýhody AM-DSB (Double Side Band) se zachovanou nosnou}:
\begin{itemize}
  \item Obě postranní pásma nesou \emph{stejnou} informaci – plýtvání spektrem.
  \item Větší část výkonu je v nosné vlně – nízká energetická účinnost.
  \item Jednoduché demodulaci (obálkový detektor).
\end{itemize}

\textbf{Varianty a jejich výhody}:
\begin{itemize}
  \item \textbf{AM-DSB-SC} (Suppressed Carrier): potlačená nosná → lepší energetická účinnost, ale složitější demodulace.
  \item \textbf{AM-SSB} (Single Side Band): přenáší jen jedno postranní pásmo → úspora spektra, ale problematické filtrování.
  \item \textbf{AM-VSB} (Vestigial Side Band): částečné potlačení jednoho pásma → kompromis (analogová TV).
\end{itemize}

% -----------------------------------------------------------
\begin{hvezdBox}
\subsection*{Úhlové modulace – PM a FM}
\end{hvezdBox}

U úhlových modulací je \textbf{obálka konstantní} $V(t) = A_c = \text{konst.}$; mění se \emph{fáze} nebo \emph{kmitočet}:

\subsubsection*{Fázová modulace (PM – Phase Modulation)}
\begin{matBox}[]
\[
  s_{PM}(t) = A_c\cos\!\bigl(\omega_c t + k_P\cdot s_M(t)\bigr)
\]
kde $k_P$ je citlivost fázového modulátoru.
\end{matBox}

\subsubsection*{Kmitočtová modulace (FM – Frequency Modulation)}
\begin{matBox}[]
\[
  s_{FM}(t) = A_c\cos\!\left(\omega_c t + k_F\int_{-\infty}^{t} s_M(\tau)\,\dd\tau\right)
\]
kde $k_F$ je citlivost kmitočtového modulátoru.
\end{matBox}

\textbf{Vztah PM a FM}: okamžitý kmitočet je derivací fáze:
\[
  \omega_i(t) = \omega_c + k_F\cdot s_M(t) \quad \text{(FM)}
  \qquad
  \omega_i(t) = \omega_c + k_P\cdot\frac{\dd s_M}{\dd t} \quad \text{(PM)}
\]
FM s integrovaným modulačním signálem = PM a naopak. Jsou proto \textbf{vzájemně zaměnitelné} s vhodným předzpracováním.

\textbf{Výhody úhlových modulací}:
\begin{itemize}
  \item Konstantní obálka → energeticky výhodné, možné použití nelineárních zesilovačů.
  \item Lepší odolnost vůči šumu (FM oproti AM).
\end{itemize}

\subsubsection*{Spektrum FM při harmonickém modulačním signálu}

Pro $s_M(t) = A_M\cos(\omega_M t)$, modulační index $\beta = k_F A_M / \omega_M$:
\[
  s_{FM}(t) = A_c\sum_{n=-\infty}^{\infty} J_n(\beta)\cos\!\bigl((\omega_c + n\omega_M)t\bigr)
\]
kde $J_n(\beta)$ jsou Besselovy funkce 1. druhu řádu $n$.

\begin{itemize}
  \item Spektrum obsahuje \emph{nekonečně mnoho} postranních pásem (na rozdíl od AM).
  \item Šířka pásma (Carsonův vztah):
    \[
      B \approx 2(\Delta f_{max} + f_{M,max}) = 2(1+\beta)f_{M,max}
    \]
  \item Kmitočtový zdvih: $\Delta f_{max} = k_F A_M / (2\pi)$.
\end{itemize}

\newpage
% ============================================================
\section{Základní charakteristiky náhodného procesu, stacionarita, ergodicita, bílý šum}
% ============================================================

\begin{hvezdBox}
\subsection*{Náhodný signál (náhodný proces)}
\end{hvezdBox}

\textbf{Náhodný proces} $\xi(t, \zeta)$ je parametrizovaný soubor náhodných jevů $\zeta \in \Omega$:
\begin{itemize}
  \item Pro \emph{fixní} $\zeta = \zeta_i$: $\xi(t, \zeta_i)$ je jedna \textbf{realizace} – deterministický signál.
  \item Pro \emph{fixní} $t = t_0$: $\xi(t_0, \zeta)$ je \textbf{náhodná veličina}.
  \item Náhodné signály mají náhodný (nikoli chaotický) průběh – popisujeme je pravděpodobnostními charakteristikami.
  \item Obsahují \emph{neznámou informaci} nebo šum.
\end{itemize}

\subsection*{Pravděpodobnostní charakteristiky}

\begin{itemize}
  \item \textbf{CDF} (Cumulative Density Function) $F_\xi(x;t) = P(\xi(t) \leq x)$: pravděpodobnost, že hodnota nepřekročí $x$.
  \item \textbf{PDF} (Probability Density Function) $f_\xi(x;t) = \dfrac{\dd F_\xi}{\dd x}$: hustota pravděpodobnosti.
  \item \textbf{Střední hodnota}: $\mu(t) = E[\xi(t)] = \int_{-\infty}^{\infty} x\,f_\xi(x;t)\,\dd x$.
  \item \textbf{Rozptyl}: $\sigma^2(t) = E[(\xi(t)-\mu(t))^2]$ – odchylka od střední hodnoty.
  \item \textbf{Autokorelační funkce}: $R_\xi(t_1, t_2) = E[\xi(t_1)\,\xi^*(t_2)]$ – závislost hodnot ve dvou časech.
\end{itemize}

% -----------------------------------------------------------
\begin{hvezdBox}
\subsection*{Stacionární náhodné signály}
\end{hvezdBox}

\textbf{Stacionární} náhodný proces = jeho pravděpodobnostní popis \textbf{nezávisí na volbě počátku časové osy}.

\begin{infoBox}[SSS – Stacionarita v užším smyslu (Strict Sense Stationary)]
Kompletní pravděpodobnostní popis (hustota $n$-tého řádu) nezávisí na časovém posunu:
\[
  f_\xi(x_1, \ldots, x_n; t_1, \ldots, t_n) = f_\xi(x_1, \ldots, x_n; t_1+\Delta t, \ldots, t_n+\Delta t)
\]
\end{infoBox}

\begin{infoBox}[WSS – Stacionarita v širším smyslu (Wide Sense Stationary)]
Pouze vybrané charakteristiky (střední hodnota a autokorelační funkce) jsou invariantní:
\begin{align*}
  &\mu = E[\xi(t)] = \text{konst.}\\
  &R_\xi(t_1, t_2) = R_\xi(t_2 - t_1) = R_\xi(\tau) \quad \text{(závisí jen na $\tau = t_2 - t_1$)}
\end{align*}
\end{infoBox}

SSS $\Rightarrow$ WSS, ale WSS $\not\Rightarrow$ SSS.

% -----------------------------------------------------------
\subsection*{Wiener-Chinčinova věta}

Pro WSS náhodný proces platí, že \textbf{spektrální hustota výkonu} $S_\xi(\omega)$ je Fourierovou transformací autokorelační funkce $R_\xi(\tau)$:

\begin{matBox}[Wiener-Chinčinova věta]
\[
  S_\xi(\omega) = \FT\{R_\xi(\tau)\} = \int_{-\infty}^{+\infty} R_\xi(\tau)\, e^{-j\omega\tau}\,\dd\tau
\]
\end{matBox}

Zpětně: $R_\xi(\tau) = \FT^{-1}\{S_\xi(\omega)\}$.

% -----------------------------------------------------------
\begin{hvezdBox}
\subsection*{Bílý šum (White Noise)}
\end{hvezdBox}

\textbf{Bílý šum} je idealizovaný model náhodného signálu s \textbf{konstantní spektrální hustotou výkonu}:

\begin{matBox}[Bílý šum]
\[
  S_\xi(\omega) = S_0 = \text{konst.} \quad (\text{pro všechna }\omega)
\]
\[
  R_\xi(\tau) = S_0\,\delta(\tau) \quad \text{(autokorelační funkce = Dirac)}
\]
\end{matBox}

\begin{itemize}
  \item Hodnoty v různých časech jsou \textbf{nekorelované} (ACF = Dirac v $\tau = 0$).
  \item \textbf{Nekonečný výkon} ($\int S_\xi\,\dd\omega = \infty$) – fyzikálně nerealizovatelný.
  \item \textbf{Nulová střední hodnota}.
  \item „Bílý" analogicky ke světlu – zastoupeny všechny frekvence stejně.
\end{itemize}

\subsubsection*{AWGN (Additive White Gaussian Noise)}

Model šumu v komunikačních systémech:
\begin{itemize}
  \item \textbf{Aditivní}: přičítá se k signálu – $r(t) = s(t) + n(t)$.
  \item \textbf{Bílý}: konstantní spektrální hustota výkonu $N_0/2$.
  \item \textbf{Gaussovský}: PDF hodnot je Gaussovo (normální) rozdělení $\mathcal{N}(0, \sigma^2)$.
  \item Nekorelovaný – vzorky v různých časech jsou statisticky nezávislé.
\end{itemize}

% -----------------------------------------------------------
\begin{hvezdBox}
\subsection*{Ergodické náhodné signály}
\end{hvezdBox}

\textbf{Ergodické signály}: třída WSS náhodných signálů, pro které \emph{časové charakteristiky} (průměrování po jedné realizaci v čase) jsou \textbf{rovny charakteristikám na množině realizací} (průměrování přes všechny realizace v daném čase).

\begin{matBox}[Ergodická střední hodnota]
\[
  \mu = \lim_{T\to\infty}\frac{1}{2T}\int_{-T}^{T}\xi(t,\zeta_i)\,\dd t \quad \text{(pro skoro všechna } \zeta_i)
\]
\end{matBox}

Klíčové vlastnosti:
\begin{itemize}
  \item Časová charakteristika je \textbf{determinovaná} (nulový rozptyl mezi realizacemi).
  \item Nepotřebujeme průměrovat přes \emph{množinu realizací} – stačí jedna dlouhá realizace.
  \item Praktický význam: měření statistik ze záznamu reálného signálu.
\end{itemize}

\subsection*{Zjednodušení výpočtu pro stacionární a ergodické signály}

\begin{center}
\begin{tabular}{lll}
\toprule
\textbf{Charakteristika} & \textbf{Obecně} & \textbf{WSS ergodický} \\
\midrule
Střední hodnota & $E[\xi(t)]$ & $\Av[\xi(t, \zeta_i)]$ (jedna realizace) \\
Výkon & $E[|\xi(t)|^2]$ & $\Av[|\xi(t)|^2]$ (časový průměr) \\
ACF & $E[\xi(t+\tau)\xi^*(t)]$ & $\Av[\xi(t+\tau)\xi^*(t)]$ (korelace v čase) \\
\bottomrule
\end{tabular}
\end{center}

\newpage
% ============================================================
\section*{Přehled klíčových vzorců a vztahů}
% ============================================================
\addcontentsline{toc}{section}{Přehled klíčových vzorců}

\begin{infoBox}[Fourierovské transformace – přehled]
\begin{align*}
  \text{FS:} & \quad c_n = \frac{1}{T_0}\int_{T_0} s(t)e^{-jn\omega_0 t}\dd t, \quad s(t) = \sum_n c_n e^{jn\omega_0 t} \\[4pt]
  \text{FT:} & \quad S(\omega) = \int_{-\infty}^{\infty} s(t)e^{-j\omega t}\dd t, \quad s(t) = \frac{1}{2\pi}\int_{-\infty}^{\infty} S(\omega)e^{j\omega t}\dd\omega \\[4pt]
  \text{DtFT:} & \quad S(e^{j\Omega}) = \sum_{k=-\infty}^{\infty} s[k]e^{-j\Omega k} \\[4pt]
  \text{DFT:} & \quad D[n] = \sum_{k=0}^{N-1} d[k]e^{-jnk\frac{2\pi}{N}}
\end{align*}
\end{infoBox}

\begin{matBox}[Konvoluční věta]
$\FT\{s_1 * s_2\} = S_1 \cdot S_2$ \hfill (konvoluce $\leftrightarrow$ násobení)
\end{matBox}

\begin{matBox}[Shannon]
$f_{sa} > 2f_m$ \hfill (vzorkovací podmínka)
\end{matBox}

\begin{matBox}[BIBO stabilita]
Spojitá: póly v levé polorovině. Diskrétní: póly uvnitř jednotkové kružnice.
\end{matBox}

\begin{matBox}[Bílý šum]
$S_\xi(\omega) = S_0$, $R_\xi(\tau) = S_0\delta(\tau)$ \hfill (Wiener-Chinčin)
\end{matBox}

\begin{matBox}[AM modulace]
$s_{AM}(t) = A_c[1 + m\cdot s_M(t)]\cos(\omega_c t)$
\end{matBox}

\begin{matBox}[FM modulace]
$s_{FM}(t) = A_c\cos\!\left(\omega_c t + k_F\int s_M\,\dd t\right)$, \quad $B \approx 2(\Delta f + f_{M,max})$
\end{matBox}

\begin{matBox}[Komplexní obálka]
$s(t) = \operatorname{Re}\{\tilde{s}(t)\cdot e^{j\omega_c t}\}$, \quad $\tilde{s}(t) = z(t)\cdot e^{-j\omega_c t}$
\end{matBox}

\end{document}
